\chapter{INTRODUCTION TO THE PROJECT}

\section{Background of the Project}

Unmanned Aerial Vehicles (UAVs) have become an essential technology 
in modern communication and monitoring systems. They are widely used 
in applications such as surveillance, disaster management, environmental 
monitoring, agriculture, and search-and-rescue operations 
\cite{ghamari2022unmanned,mohsan2023unmanned}. 
In many real-world missions, multiple UAVs operate together to perform 
coordinated tasks. These UAVs communicate through wireless networks 
known as UAV Ad Hoc Networks (UAANETs) \cite{rovira2022review,nawaz2021uav}.

Communication between UAVs and Ground Control Stations is commonly 
achieved using the MAVLink protocol. MAVLink is a lightweight 
communication protocol designed for transmitting telemetry data, 
navigation commands, and control messages \cite{koubaa2019micro}. 
Although MAVLink is efficient and widely adopted, it relies on fixed 
transmission rates that do not adapt to changing network conditions 
\cite{cordill2025comprehensive}.

As the number of UAVs increases, communication traffic also increases. 
This leads to network congestion, packet loss, and increased 
transmission delays \cite{anwekar2023analysis}. 
Loss of critical messages such as control commands may affect mission 
safety and reliability. Therefore, improving communication efficiency 
and reliability in UAV networks is an important research problem 
\cite{verma2024adaptive}.

To address these challenges, this project proposes a Machine 
Learning-based message prioritization and congestion-aware rate 
control system. The proposed system uses the Random Forest 
algorithm to classify MAVLink messages into priority levels. 
Based on congestion levels, transmission rates are adjusted 
dynamically to reduce packet loss and improve communication 
reliability.

The final framework extends this idea by adding reinforcement 
learning-based scheduling and aggregation on top of the ML-based 
prioritization and congestion control layers. The learned policy 
selects which queue to serve and how many packets to aggregate in 
each transmission cycle, making the communication process adaptive 
to both traffic urgency and channel conditions.

\section{Objectives and Contributions}

The main objectives of this project are to:

\begin{itemize}
	\item classify MAVLink messages according to operational priority using a Machine Learning model,
	\item regulate transmission rates based on congestion conditions,
	\item jointly optimize message scheduling and packet aggregation using a learned reinforcement learning policy,
	\item maintain backward compatibility with the MAVLink protocol, and
	\item evaluate the combined framework against a paper-style baseline using throughput and packet-loss metrics.
\end{itemize}

The major contributions are:

\begin{enumerate}
	\item a semantic message prioritization module based on Random Forest classification,
	\item a heuristic congestion-aware rate control mechanism,
	\item a closed-loop RL-based scheduling and aggregation module,
	\item scenario-wise comparison against the chosen baseline paper, and
	\item a lightweight implementation suitable for edge-style UAV communication settings.
\end{enumerate}

\section{Organization of the Paper}

The rest of the paper is organized as follows. Chapter 2 reviews the 
related work and identifies the baseline paper used for comparison. 
Chapter 3 presents the design and mathematical formulation of the 
proposed framework. Chapter 4 describes the implementation details, 
including the ML model and RL training setup. Chapter 5 discusses the 
experimental results and the comparison with the baseline. Chapter 6 
concludes the work and outlines future research directions.

\section{Problem Statement}

In MAVLink-based UAV communication systems, multiple types of 
messages such as command messages, telemetry data, navigation 
updates, and sensor information are transmitted simultaneously 
\cite{koubaa2019micro}. These messages are usually transmitted 
at fixed rates without considering current network conditions.

When the incoming message rate exceeds the available network 
bandwidth, packets accumulate in transmission queues. If the 
queue capacity is exceeded, packets are dropped, resulting in 
increased packet loss and reduced Quality of Service (QoS) 
\cite{anwekar2023analysis}. This problem becomes more critical 
when high-priority messages are delayed or lost.

Existing communication systems mainly rely on simple scheduling 
techniques and static transmission rates, which are not suitable 
for highly dynamic UAV environments \cite{li2018prioritizing,
halder2022dynamic}. There is a lack of intelligent systems that 
can dynamically adjust message transmission rates based on 
congestion and message importance \cite{verma2024adaptive}.

Therefore, there is a need to develop an adaptive communication 
system that can prioritize important messages using Machine 
Learning techniques and regulate transmission rates based on 
congestion levels. The proposed project addresses this challenge 
by integrating Random Forest-based message classification with 
congestion-aware rate control to improve overall communication 
performance in UAV networks.